\section[Médiation et autonomie]{La médiation numérique et l'autonomie des chercheurs}
Le passage d'un outil figé à une plateforme dynamique ne relève pas uniqueemnt d'un alignement technique, il s'inscrit aussi dans une mutation des pratiques professionnelles. À la bibliothèque Cujas, cette conduite du changement doit permettre d'alléger la charge de travail et valoriser l'expertise métier des agents.

\subsection{La médiation et l'autonomie des chercheurs}
La mise en place d'un back-office collaboratif redistribue les rôles et modifie le quotidien des agents. L'analyse fonctionnelle réalisée dans le cadre de la réalisation de l'audit permet d'identifier l'évolution du rôle d'accompagne des personnels des bibliothèques. 

Les bibliothécaires, étant au contact du public, ont besoin de connaître avec précision les services de la bibliothèque et d'être en mesure d'expliquer son fonctionnement. Ils doivent aussi pouvoir répondre aux questions éventuelles que les usagers pourraient être amenés à lui poser. Sa motivation première est d'accompagner les usagers pour les rendre autonome dans leur appropriation des outils. Face à l'arrivée de briques technologiques complexes, comme la visionneuse Mirador ou le serveur IIIF, la conduite du changement place le bibliothécaire dans un rôle de médiateur permettant de rendre autonome le chercheur. Dans cette nouvelle posture, il est amené à rédiger des tutoriels ou des pages d'aide pour les différents publics de la bibliothèque.

Le responsable de la bibliothèque numérique se doit d'avoir une bonne compréhension globale de son infrastructure. C'est sur cette connaissance qu'il pourrait s'appuyer dans le cadre de la réalisation de projet scientifique en collaboration avec les chercheurs afin de pouvoir mettre en avant les collections numérisées.La future bibliothèque numérique doit devenir un espace collaboratif favorisant l'appropriation active et autonome des sources présentes à l'intérieur.

\subsection{La co-conception des programmes documentaires}
L'histoire de CujasNum est marquée par la collaboration scientifique, notamment pour la liste Pfister-Roumy. Pour que la future plateforme demeure une référence, la bibliothèque peut envisager d'ouvrir la conception de la politique documentaire à un comité scientifique de chercheurs. Ce comité aura pour mission de définir et de prioriser de manière collaborative les futurs chantiers de numérisation. Par un exemple, un des enseignants-chercheurs proposait et d'initier en ce sens un nouveau projet de numérisation, nommé la \enquote{liste Pfister-Roumy 2.0} ciblant \enquote{les manuels de droits civils et les grands revues de la période de 1800 à 1940}. Cette collaboration permettra à la bibliothèque Cujas de constituer des projets scientifiques cohérents et pertinents pour les chercheurs. De cette façon, le chercheur prend une part active à la conception de la future bibliothèque numérique.

\subsection{L'éditorialisation participative}
L'un des avantages de cette refonte est l'intégration des chercheurs et des étudiants avancés au sein du processus d'éditorialisation et de valorisation des collections. CujasNum laissait peu de place à l'éditorialisation, car le logiciel ne le permettait pas. La nouvelle plateforme doit pouvoir proposer des espaces d'éditorialisation participative. Ce moyen de fonctionnement permet ainsi de s'appuyer sur un cercle vertueux d'insertion des jeunes chercheurs.

Elina Leblanc met en avant ce concept \enquote{d'éditorialisation comme co-création scientifique}. Elle analyse le passage de l'usager passif au contributeur actif : en offrant aux chercheurs ou aux doctorants des rôles de \enquote{Contributeurs} pour rédiger des focus ou annoter des documents IIIF, ainsi la bibliothèque crée un espace de co-construction de nouveaux savoirs. Toutefois, la responsabilité de la bibliothèque demeure inchangée, elle doit \enquote{encadrer techniquement et scientifiquement ce travail collaboratif en back-office pour en garantir la qualité et la fiabilité sémantique}. En ouvrant ses outils éditoriaux à la communauté académique, la bibliothèque numérique s'affirme comme un espace d'apprentissage et de diffusion scientifique partagé.
\bigskip

La refonte de la bibliothèque numérique de Cujas démontre trois enjeux majeurs qui sont humains, financiers et organisationnels. Elle amène à réfléchir à un décloisonnement des pratiques internes, envisageant de nouveaux circuits de travail. Le rôle des bibliothèques numériques n'est plus de simplement mettre en ligne les documents numérisés, mais d'impliquer directement la communauté scientifique dedans. Les bibliothèques doivent être en mesure de les accompagner vers une réelle autonomie de traitements des données et des outils proposés. Il est donc nécessaire d'accompagner les bibliothécaires eux-mêmes dans cette conduite du changement, afin de permettre à cette bibliothèque numérique de devenir un véritable espace de vie intellectuelle, de médiation culturelle et de partage scientifique du savoir.